\documentclass[10pt,twocolumn]{article}
\usepackage[margin=0.62in]{geometry}
\usepackage{graphicx,amsmath,booktabs,hyperref,enumitem}
\title{TABasco: Task-Aware Browser Tab Optimization using Reinforcement Learning}
\author{Team members: \textbf{[Insert full names and roll numbers]}\\GitHub: \url{https://github.com/Aarsussh/TABasco}}
\date{CSL3090 Artificial Intelligence -- Assignment 2 (October 2026)}
\begin{document}\maketitle
\begin{abstract}\small
TABasco studies browser tab management as a finite-horizon Markov decision process (MDP) under active-memory pressure and changing tasks. We implement a synthetic browser session, a reward-driven Deep Q-Network (DQN), and LRU, LFU, and random evaluation baselines. On 60 held-out session seeds, the trained DQN obtained mean return $-48.73$ (SD $23.51$), compared with $-59.40$ (SD $5.94$) for LRU, $-59.36$ (SD $5.83$) for LFU, and $-59.13$ (SD $5.88$) for Random. The DQN's invalid-action rate was $0.448$, substantially above baseline rates, and its return variability was large. These results show a functioning RL experiment, not reliable policy superiority.
\end{abstract}
\section{Problem and MDP}
Browsers keep many tabs available while active pages consume finite memory. Suspending a tab saves active memory but may interrupt a task and require restoration; closing it risks lost work and recovery effort. TABasco models this sequential tradeoff without controlling a real browser.

The state $s_t$ contains $N=12$ tab slots and five global features. Each slot has eight normalized values: occupied, active, memory size, time since visit, visit frequency, current-task match, estimated revisit probability, and unsaved-work risk. Global values encode active-memory/budget, current task, remaining horizon, elapsed fraction, and occupancy. The observation is a fixed 101-dimensional vector. Tasks range over four contexts. Tab memory sizes are sampled uniformly from 0.45 to 1.35 synthetic units. Revisit probability decays with age, increases with frequency, and is scaled by task affinity. Sessions also include task switches, stochastic tab arrivals, aging/expiration, revisits, and closed-tab history recovery.

At each step the agent chooses one of three operations on one slot: KEEP, SUSPEND, or CLOSE. Action $a=3i+o$ denotes slot $i$ and operation $o$. One action advances the simulated session by one time step. Empty/out-of-range actions are invalid, preserve the tab slots, and incur a fixed penalty. Active tabs consume their sampled memory; suspended tabs consume zero active memory while their metadata remains available. Budget excess is measured on active memory.

The reward is the negative sum
\begin{align*}
r_t=-[&1.0\max(0,M_t-B)+C^{\text{restore}}_t+1.2C^{\text{lost}}_t\\
&+0.45C^{\text{disrupt}}_t+C^{\text{action}}_t].
\end{align*}
Restoration includes a cost for a suspended-tab return or a higher closed-tab recovery cost. Closing incurs lost-work risk. Task disruption is proportional to the tab's estimated revisit probability and whether it matches the active task. Ordinary operations cost 0.035; invalid actions cost 0.12. No reward is paid merely for reducing tab count.

\section{Learning and experiment}
The DQN is a 101-input MLP with two 128-unit ReLU layers and 36 Q outputs. It uses epsilon-greedy exploration, a uniform replay buffer, discount factor 0.97, Adam learning rate 0.001, batch size 64, Huber TD loss, gradient clipping, and a target network copied every 500 updates. Training runs on CPU for 250 episodes of at most 120 steps, seed 3090. The replay learner updates from observed transitions and reward; no baseline supplies labels. Checkpoint: \texttt{models/tabasco\_dqn.pt}.

For comparison, LRU suspends the active tab with the oldest last-visit step, LFU suspends the active tab with the lowest access count (ties by recency), and Random uniformly chooses an active tab to suspend. If no active tab is eligible, the policies choose the interface's fallback action. All four policies use identical episode seeds 90000--90059 and the same environment parameters. The DQN acts greedily. Each result represents one policy rollout per seed; variability is across those 60 episodes.

\section{Results}
\begin{table}[h]\centering\small
\caption{Held-out evaluation (60 episodes). Mean and standard deviation.}
\begin{tabular}{lrrr}\toprule
Policy & Return & Memory & Invalid\\\midrule
DQN & $-48.73\pm23.51$ & 6.28 & 0.448\\
LRU & $-59.40\pm5.94$ & 1.25 & 0.006\\
LFU & $-59.36\pm5.83$ & 1.30 & 0.006\\
Random & $-59.13\pm5.88$ & 1.36 & 0.006\\\bottomrule
\end{tabular}\end{table}

The DQN has a numerically higher mean return, but it also exhibits much greater dispersion, mean peak memory 9.45 against budget 8, and 23.23 violating steps per episode on average. LRU/LFU/Random maintain lower mean active memory (about 1.25--1.36) and lower violation count (about 0.6 steps), while paying more restoration cost (about 19.5 versus 6.9 for DQN). Useful revisit success is 0.803 for DQN and 0.876 for each baseline. DQN's invalid actions frequently target empty slots because action masking is not used. Thus the measured return difference alone is not evidence of a robustly better memory policy. Reward scale includes explicit recovery, disruption, and memory terms; results should be read with these components.

\begin{figure}[h]\centering\includegraphics[width=\columnwidth]{../results/training_returns.png}\caption{DQN training episode returns and rolling mean (window derived from log length).}\end{figure}
\begin{figure}[h]\centering\includegraphics[width=\columnwidth]{../results/return_comparison.png}\caption{Mean held-out return and one standard deviation across 60 common seeds.}\end{figure}
\begin{figure}[h]\centering\includegraphics[width=\columnwidth]{../results/memory_comparison.png}\caption{Mean active-memory usage; dashed line is the 8-unit budget.}\end{figure}

\section{Conclusion and limitations}
This implementation provides a complete, reproducible RL formulation and a working DQN experiment. The synthetic session model is an abstraction, and one 250-episode training seed is insufficient to support generalization claims. Unmasked actions allow frequent invalid choices, while baseline policies aggressively suspend tabs and incur larger restoration costs. Next steps are action masking in both exploration and greedy selection, multiple independent training seeds, sensitivity analysis of reward coefficients, and calibration from browser workload traces. Team identifiers should be completed before submission.

\clearpage\onecolumn\begin{thebibliography}{9}
\bibitem{mnih} V. Mnih et al., ``Human-level control through deep reinforcement learning,'' \emph{Nature}, vol. 518, pp. 529--533, 2015.
\bibitem{sutton} R. S. Sutton and A. G. Barto, \emph{Reinforcement Learning: An Introduction}, 2nd ed. MIT Press, 2018.
\end{thebibliography}
\end{document}
